\documentclass[../../main.tex]{subfiles}
\begin{document}

{\bf [解]}
題意より
\begin{align}
 a,b\in\mathbb{Z}_{>0} \label{eq:1}
\end{align}
である．

\paragraph{(1)}
題意より$c=a+b,\ d=a^2-ab+b^2$である．\cref{eq:1}より$c>0, d=(a-b)^2+ab>0$だから，
\begin{align}
 c^2-d=3ab>0 \qquad 4d-c^2=3(a-b)^2\ge 0 
\end{align}
だから
\begin{align}
 d < c^2 \le 4d \\
\therefore
1<\dfrac{c^2}{d}\le 4
\end{align}
であり，題意は示された．右側の等号成立条件は$a=b$の時である．


\paragraph{(2)}
$A=a^3+b^3$とおいて$A=p^n$（$p\in\text{prime}$，$n\in\mathbb{Z}$）となる条件をしらべる．
まず\cref{eq:1}から$a,b \ge 1$なので$A\in\mathbb{Z}_{\ge2}$である．
従って $n\in\mathbb{N}$として良い．

さて，$A$を因数分解すると
\begin{align}
 A = (a+b)(a^2-ab+b^2) = cd
\end{align}
である．従って$c=p^\alpha,\ d=p^\beta$（$\alpha,\beta\in\mathbb{Z}_{\ge0}$）とおける．
(1)の結果から
\begin{align}
  1<p^{2\alpha-\beta}\le4 \label{eq:2}
\end{align}
であるが，$p$が素数だから\cref{eq:2}を満たす組は
\begin{align}
  (p,2\alpha-\beta)=(2,1), (2,2), (3,1)\ \label{eq:3}
\end{align}
である．

\subparagraph{$(p,2\alpha-\beta)=(2,1)$の時}

この時$a,b$は
\begin{align}
 a+b=2^\alpha, \qquad a^2-ab+b^2=2^{2\alpha-1}
\end{align}
と書けるから$(a+b)^2=a^2-ab+b^2+3ab$に代入して
\begin{align}
 2^{2\alpha} = 2^{2\alpha-1} + 3ab \\
\therefore
 3ab = 2^{2\alpha-1}
\end{align}
となる．ここで$a,b$は自然数，$\beta$は$0$以上の整数だから，右辺が$3$を因数に持たずに不適である．

\subparagraph{$(p,2\alpha-\beta)=(2,2)$の時}

この時$a,b$は
\begin{align}
  a+b=2^\alpha, \qquad a^2-ab+b^2=2^{2\alpha-2}
\end{align}
と書ける．まず\cref{eq:2}で等号が成立するので(1)から$a=b$である．これを代入して
\begin{align}
  a=b=2^{\alpha-1}
\end{align}
となる．$a,b$は自然数だから$\alpha\ge 1$であることが必要で，よって非負の整数$n$を用いて
\begin{align}
 (a,b) = (2^n,2^n) \label{eq:4}
\end{align}
と書ける．逆にこの時
\begin{align}
 A = 2\cdot 2^{3n} = 2^{3n+1}
\end{align}
とかけて題意を満たすから，\cref{eq:4}は解である．

\subparagraph{$(p,2\alpha-\beta)=(3,1)$の時}

この時$a,b$は
\begin{align}
  a+b=3^\alpha, \qquad a^2-ab+b^2=3^{2\alpha-1} \quad\cdots\text{③} \label{eq:5}
\end{align}
である．従って第一式の二乗と第二式の三倍を計算すると
\begin{align}
 (a+b)^2 = 3^{2\alpha} \\
 3(a^2-ab+b^2) = 3^{2\alpha}
\end{align}
だから，$\alpha$を消去して
\begin{align}
 3(a^2-ab+b^2) = (a+b)^2 \\
\therefore
 (2a-b)(a-2b) = 0 \\
\therefore
  b=2a, a=2b
\end{align}
である．ここで対称性から$b=2a$の場合のみ考えよう．
\cref{eq:5}に代入して
\begin{align}
  (a,b)=\left(3^{\alpha-1},2\cdot 3^{\alpha-1}\right) 
\end{align}
である．$a,b$は自然数だから$\alpha\ge 1$なので，非負の整数$n$を用いて
\begin{align}
  (a,b)=\left(3^n,2\cdot 3^n\right) 
\end{align}
と書ける．逆にこの時
\begin{align}
  A = 9\cdot 3^{3n} = 3^{3n+2}
\end{align}
となって条件を満たす．逆の$a=2b$のパターンも同様だから
\begin{align}
  (a,b)=\left(3^n,2\cdot 3^n\right), \left(2\cdot3^n,3^n\right) \label{eq:6}
\end{align}
が解である．
\casesend


以上の場合わけで全ての場合が尽くされた．よって求める解は\cref{eq:4,eq:6}から
\begin{align}
  (a,b)=(2^n,2^n),\ (2\cdot3^n,3^n),\ (3^n,2\cdot3^n)\quad(n\in\mathbb{Z}_{\ge0})
\end{align}
である．


\newpage

\end{document}
